\documentclass[10pt]{article}
\usepackage[utf8]{inputenc}
\usepackage[T1]{fontenc}
\usepackage{amsmath}
\usepackage{amsfonts}
\usepackage{amssymb}
\usepackage[version=4]{mhchem}
\usepackage{stmaryrd}
\usepackage{bbold}

\begin{document}
$$
\angle 5
$$

1

L5: Divichlet's Theorem and applications of the FTA\\
Theorem 1.21 (Dirichlet's Theorem on primes in AP): Let $a, b \in \mathbb{Z}$ with $a, b>0$ and $(a, b)=1$. Then the arithmetic progression (AP)\\
$a, a+b, a+2 b, \ldots, a+n$,\\
contains infinitely many primes.

\begin{itemize}
  \item Unfortunately, the proof of Dirichlet's Theorem is outside the scope of this class.\\
Remark: Setting $a=b=I$ recovers the result we proved about the infinitude of prime numbers.
  \item We can use the FTA to prove special cases of Dirichlet's Theorem e.g. $a=3, b=4$.\\
Proposition 1.22: There are infinitely many primes of the form $4 n+3$ where $n$ is a non-negative integer.
  \item First, we need the following Lemma.
\end{itemize}

Lemma).23: Let $a, b \in \mathbb{Z}$. If $a$ and $b$ are expressible in the form $4 n+1$, where $n$ is an integer, then ab is also expressible in that form.\\
Proof: Let $a=4 n_{1}+1$ and $b=4 n_{2}+1$ with $n_{1}, n_{2} \in \mathbb{Z}$. Then

$$
\begin{aligned}
a b & =\left(4 n_{1}+1\right)\left(4 n_{2}+1\right) \\
& =16 n_{1} n_{2}+4 n_{1}+4 n_{2}+1 \\
& =4\left(4 n_{1} n_{2}+n_{1}+n_{2}\right)+1 \\
& =4 n+1,
\end{aligned}
$$

where $n=4 n_{1} n_{2}+n_{1}+n_{2} \in \mathbb{Z}$.\\
Proof (of Prop.1.22): Assume, by way of contradiction, that there are finitely many prime numbers of the form $4 n+3$ where $n$ is a non-negative integer, say

$$
P_{0}:=3, P_{1}, P_{2}, \cdots, P_{r} .
$$

Here, the prime numbers $P_{1}, \cdots, P_{r}$ are assumed to be distinct from $P_{0}=3$. Consider the\\
integer $N=4 P_{1} \cdots P_{r}+3$. The prime factorisation of $N$ must contain a prime number $p$ of the desired form or it would otherwise contain only prime numbers of the form $4 n+1$ where $n$ is a positive integer, and hence be of the form $4 n+1$ where $n$ is a positive integer by Lemmal.23.\\
Lo $p=p_{0}=3$ or $p=p_{i}$ for some $i=1,2, \ldots, r$.\\
Case I $\left(p=p_{0}=3\right)$ : Then 31 N implies that $3 \mid N-3$ by Proposition 1.2, and we have that $314 P_{1} P_{2} \cdots P_{r}$. By Corollary 1.15, we now have that 314 or 31 Pi for some $i=1,2, \ldots, r$. Contradiction I

Case II ( $p=p_{i}$ for some $i=1,2, \ldots, r$ ): Then $P_{i} \mid N$ implies that $P_{i} \mid N-4 P_{1} P_{2} \cdots P_{r}$ by Proposition 1.2, and we have that $p_{i} \mid 3$, a contradiction! Thus there are infinitely\\
many prime numbers of the desired form.\\
Remark: The proof of Proposition $1.22 \mathrm{can}^{\prime}$ 't be extended to prove Dirichlet's Theorem in full generality i.e. Lemma 1.23 does not hold for integers of the form $4 n+3$ (exercise for home).

\section*{Congruences}
\begin{itemize}
  \item Up to the beginning of the $19^{\text {th }}$ century number theory consisted of brilliant results but with no connecting theme.
  \item In 1801, Gauss in "Disquisitiones Arithmeticae" unified number theory $\rightarrow$ and developed the theory of congruences. Gauss was only 24 years of age!\\
Definition: Let $a, b, m \in \mathbb{Z}$ with $m>0$. Then $a$ is said to be congrnent to $b \operatorname{modulo} m$, denoted $a \equiv b(\bmod m)$, if $m \mid a-b$. If\\
$a \equiv b(\bmod m)$, then $m$ is said to be the modulus of the congruence. The notation $a \neq b(\bmod m)$ means $a$ is not congment to $b$ modulo $m$ i.e. $m+a-b$.\\
Remark: The modulus in a congruence is defined to be a positive integer. dometimes this is assumed implicitly.\\
Example: $\cdot 4 \mid 25-1$, so $25 \equiv 1(\bmod 4)$
  \item I la-b for all $a, b \in \mathbb{Z}$, so $a \equiv b(\bmod \pm)$ for all $a, b \in \mathbb{Z}$.
  \item $5 \not+-7-2$, so $-7 \not \equiv 2(\bmod 5)$
\end{itemize}

Proposition 2.1: Congruence modulo $m$ is an equivalence relation on $\mathbb{Z}$.\\
Proof: Let $a \in \mathbb{Z}$. Since any integer divides 0 , we have $m / O$, from which $m / a-a$ and $a \equiv a(\bmod m)$. Thus Congruence modulo $m$ is reflexive on $\mathbb{Z}$. Let $a, b \in \mathbb{Z}$ and assume\\
that $a \equiv b(\bmod m)$. Then $m \mid a-b$, and (6) so we have $m \mid(-1)(a-b)$, or equivalently, $m(b-a$. Thus $b \equiv a(\bmod m)$, and congruence modulo $m$ is symmetric. Let $a, b, c \in \mathbb{Z}$ and assume that $a \equiv b(\operatorname{modm})$ and $b \equiv c(\bmod m)$. Then $m \mid a-b$ and $m \mid b-c$, and so $m /(a-b)+(b-c)$ by Proposition 1.2. Thus $m \mid a-c$, and so $a \equiv c(\bmod m)$. Thus congruence modulo $m$ is transitive on $\mathbb{Z}$. We have proved that congruence modulom is an equivalence velation on $\mathbb{Z}$.

Consequence 2.2: $\mathbb{Z}$ is partitioned into equivalence classes under congruence modulo $m$.

\begin{itemize}
  \item For a given $x \in \mathbb{Z}$, let $[x]$ denote the equivalence class of $x$ modulo $m$. This should not be confused with the floor function.
\end{itemize}

\section*{Example: The equivalence classes of $\mathbb{Z}$}
under congruence modulo 4 are

$$
\begin{aligned}
{[0] } & =\{x \in \mathbb{Z}: x \equiv 0(\bmod 4)\} \\
& =\{x \in \mathbb{Z}: 4 \mid x-0\} \\
& =\{x \in \mathbb{Z}: 4 \mid x\} \\
& =\{x: x=4 n \text { for some } n \in \mathbb{Z}\} \\
& =\{\ldots,-8,-4,0,4,8, \ldots\},
\end{aligned}
$$

and similarly,\\
$[1]=\{x: x=4 n+1$ for some $n \in \mathbb{Z}\}$

$$
=\{\ldots,-7,-3,1,5,9, \ldots\}
$$

$[2]=\{x: x=4 n+2$ for some $n \in \mathbb{Z}\}$.

$$
=\{\ldots,-6,-2,2,6,10, \ldots\}
$$

$[3]=\{x: x=4 n+3$ for some $n \in \mathbb{Z}\}$

$$
=\{\ldots,-5,-1,3,7,11, \ldots\}
$$

$\mathbb{Z}$ is partitioned into four equivalence classes\\
under congruence modulo 4 .

\begin{itemize}
  \item From this example we see that any integer is congruent to $0,1,2$, or 3 modulo 4 .\\
Definition: A set of integers such that every integer is congrnent modulo $m$ to exactly one integer of the set is said to be a Complete residue system modulo $m$.\\
Example: - $\{0,1,2,3\}$ is a complete residue system modulo 4.
  \item $\{6,-11,19,1988\}$ is also a complete residue Syrtem modulo 4.\\
i.e. $6 \equiv 2(\bmod 4),-11 \equiv 1(\bmod 4), 19 \equiv 3(\bmod 4)$, $1988 \equiv 0(\bmod 4)$.\\
Proposition 2.3 : The set $\{0,1,2, \ldots, m-1\}$ is a complete residue system modulo $m$.\\
Proof: We first show that every integer is congruent modulo $m$ to at least one\\
integer in $\{0,1,2, \ldots, m-1\}$. Let $a \in \mathbb{Z}$. (9) Then by the Dividion algorithm, there exist $q, r \in \mathbb{Z}$ such that
\end{itemize}

$$
a=m q+r, \quad 0 \leq r<m .
$$

Now $a=m q+r$ implies that $a-r=m q$, so $m \mid a-r$, or equivalently, $a \equiv r(\bmod m)$.\\
Furthermore, $0 \leq r<m$ implies that $r \in\{0,1,2, \ldots, m-1\}$, which shows the initial claim.

We now show that every integer is congruent modulo $m$ to at most one integer in $\{0,1,2, \cdots, m-1\}$, which would complete the proof. Let $a \in \mathbb{Z}$ and $r_{1}, r_{2} \in\{0,1, \ldots, m-1\}$ such that $a \equiv r_{1}(\bmod m)$ and $a \equiv r_{2}(\bmod m)$.\\
Then $r_{1} \equiv r_{2}(\bmod m)$, or equivalently $m \mid r_{1}-r_{2}$. Since $r_{1}, r_{2} \in\{0,1, \ldots, m-1\}$, we have that

$$
-(m-1) \leq r_{1}-r_{2} \leq m-1 .
$$

Since $m \mid r_{1}-r_{2}$, we have that $r_{1}-r_{2}=0$ (10) i.e. $r_{1}=r_{2}$. Thus every integer is congruent modulo $m$ to at most one integer in\\
$\{0,1,2, \ldots, m-1\}$.\\
Definition: The set $\{0,1,2, \ldots, m-1\}$ is said to be the set of least non-negative residues modulo $m$.

Proposition 2.4 : Let $a, b, c, d \in \mathbb{Z}$ such that $a \equiv b(\bmod m)$ and $c \equiv d(\bmod m)$. Then\\
(a) $a+c \equiv b+d(\operatorname{modm})$\\
(b) $\quad a c \equiv b d(\bmod m)$

Proof: Since $a \equiv b(\bmod m)$ and $c \equiv d(\bmod m)$, we have that $m \mid a-b$ and $m \mid c-d$. For (a), note that $m \mid(a-b)+(c-d)$, which is equivalent to $m \mid(a+c)-(b+d)$, from which we have

$$
a+c \equiv b+d(\bmod m) .
$$

For (b), note that $m / a-b$ implies that\\
$m \mid c(a-b)$, and that $m \mid c-d$ implies that (11) $m \mid b(c-d)$. Thus $m \mid c(a-b)+b(c-d)$, which is equivalent to $m \mid a c-b d$, and so

$$
a c \equiv b d(\bmod m) .
$$


\end{document}